%%%%%%%%%%%%%%%%%%%%%%%%%%%%%%%%%%%%%%%%%
% Beamer Presentation
% LaTeX Template
% Version 1.0 (10/11/12)
%
% This template has been downloaded from:
% http://www.LaTeXTemplates.com
%
% License:
% CC BY-NC-SA 3.0 (http://creativecommons.org/licenses/by-nc-sa/3.0/)
%
%%%%%%%%%%%%%%%%%%%%%%%%%%%%%%%%%%%%%%%%%

%----------------------------------------------------------------------------------------
%	PACKAGES AND THEMES
%----------------------------------------------------------------------------------------

\documentclass[aspectratio=169]{beamer} %im 16 zu 9 Format
%\documentclass{beamer} %im 4 zu 3 Format

\mode<presentation> {

% The Beamer class comes with a number of default slide themes
% which change the colors and layouts of slides. Below this is a list
% of all the themes, uncomment each in turn to see what they look like.

%\usetheme{default}
%\usetheme{AnnArbor}
%\usetheme{Antibes}
%\usetheme{Bergen}
%\usetheme{Berkeley}
%\usetheme{Berlin}       %gefällt mir vom Stil gut, ist aber nicht klassisch
%\usetheme{Boadilla}
%\usetheme{CambridgeUS}
%\usetheme{Copenhagen}
%\usetheme{Darmstadt}
%\usetheme{Dresden}
%\usetheme{Frankfurt}
%\usetheme{Goettingen}
%\usetheme{Hannover}
%\usetheme{Ilmenau}
%\usetheme{JuanLesPins}
%\usetheme{Luebeck}
\usetheme{Madrid}     %klassischer Stil, aber nichts besonderes
%\usetheme{Malmoe}
%\usetheme{Marburg}
%\usetheme{Montpellier}
%\usetheme{PaloAlto}
%\usetheme{Pittsburgh}
%\usetheme{Rochester}
%\usetheme{Singapore}
%\usetheme{Szeged}
%\usetheme{Warsaw}

% As well as themes, the Beamer class has a number of color themes
% for any slide theme. Uncomment each of these in turn to see how it
% changes the colors of your current slide theme.

%\usecolortheme{albatross}
%\usecolortheme{beaver}
%\usecolortheme{beetle}
%\usecolortheme{crane}
%\usecolortheme{dolphin}
%\usecolortheme{dove}
%\usecolortheme{fly}
%\usecolortheme{lily}
%\usecolortheme{orchid}
%\usecolortheme{rose}
%\usecolortheme{seagull}
%\usecolortheme{seahorse}
%\usecolortheme{whale}
%\usecolortheme{wolverine}

%\setbeamertemplate{footline} % To remove the footer line in all slides uncomment this line
%\setbeamertemplate{footline}[page number] % To replace the footer line in all slides with a simple slide count uncomment this line

%\setbeamertemplate{navigation symbols}{} % To remove the navigation symbols from the bottom of all slides uncomment this line
}

\usepackage{graphicx} % Allows including images
\usepackage{booktabs} % Allows the use of \toprule, \midrule and \bottomrule in tables
\usepackage{pdfpcnotes}
\usepackage{listings}
\usepackage[ngerman]{babel}
\usepackage[ngerman]{datetime}
\usepackage[T1]{fontenc} %Damit Umlaute auch beim Kopieren richtig dargestellt werden
\usepackage{csquotes}

\usepackage{color}

\usepackage{pgfgantt} %GANT Diagramm

% Farbdefinitionen
\definecolor{myblue}{RGB}{0, 51, 102}
\definecolor{myorange}{RGB}{255, 102, 0}

\usepackage{tikz}
\usetikzlibrary{shapes.arrows}

\tikzset{
    myarrow/.style={
        draw,
        fill=orange,
        single arrow,
        minimum height=9ex,
        single arrow head extend=1ex
    }
}

\newcommand{\arrowright}{%
\tikz [baseline=-0.5ex]{\node [myarrow,rotate=0] {};}
}
\newcommand{\arrowleft}{%
\tikz [baseline=-0.5ex]{\node [myarrow,rotate=180] {};}
}

%extra blocks

\newenvironment<>{extrablock}[1]{%
\setbeamercolor{block title}{bg=orange!90!black}%
  \begin{block}#2{#1}}{\end{block}}

\renewcommand{\figurename}{Abbildung}

\lstset{
  basicstyle=\ttfamily,
  columns=fullflexible,
  escapeinside = {(*@}{@*)},
  numbers=left,
  xleftmargin=2.5em,
  frame=single,
  breaklines=true,
}

%Für Zeilenänderungen
\let\origthelstnumber\thelstnumber
\makeatletter
\renewcommand\thelstnumber{%
    \ifnum\value{lstnumber}>0
        \origthelstnumber
    \else
        \ifnum\value{lstnumber}=-1
            \ldots
        \fi
    \fi
}

% for starting line numbers
\newcommand*\startnumber[1]{%
    \setcounter{lstnumber}{\numexpr#1-1\relax}%
}

% for stopping line numbers
\newcommand*\stopnumber{%
    \startnumber{-2}%
}
\makeatother

%----------------------------------------------------------------------------------------
%	TITLE PAGE
%----------------------------------------------------------------------------------------

\title[Fachpräsentation]{Nutzung von Daten ohne vorgegebene Zielvariable (unlabelled data) durch semiüberwachtes Lernen} % The short title appears at the bottom of every slide, the full title is only on the title page

\author{Jasper Bremenkamp} % Your name
\institute[IU] % Your institution as it will appear on the bottom of every slide, may be shorthand to save space
{
IU Internationale Hochschule \\ % Your institution for the title page
\medskip
\textit{jasper.bremenkamp@iu-study.org}\\ % Your email address
https://jasper.bremenkamp.de
}
\date{\today} % Date, can be changed to a custom date

\begin{document}

\begin{frame}
    \titlepage % Print the title page as the first slide
    
    \pnote{Herzlich willkommen zu [...] }
    \pnote{[Thema, Datum und eigenen Namen nennen]}
\end{frame}

\begin{frame}
    \frametitle{Agenda} % Table of contents slide, comment this block out to remove it
    \tableofcontents % Throughout your presentation, if you choose to use \section{} and \subsection{} commands, these will automatically be printed on this slide as an overview of your presentation

    \pnote{0-5}
    
    \pnote{0. Kommen wir direkt zur Agenda der heutigen Präsentation!}
    \pnote{1. Einleitung in das Thema, semiüberwachtes Lernen auf den ersten Blick, dessen Bedeutung und das inhaltliche Ziel der Präsentation}
    \pnote{2. Nach der Erstbetrachtung gehen wir auf \dots ein}

    \pnote{3. dann beleuchten wir \dots}

    \pnote{4. Wir schauen, was wir für Szenarien haben, wo wir semiüberwachtes Lernen nutzen können}

    \pnote{5. Zum Abschluss fasse ich das gezeigte zusammen und wir erläutern die Vorteile, aber auch Grenzen von semiüberwachtem Lernen}
\end{frame}

%----------------------------------------------------------------------------------------
%	PRESENTATION SLIDES
%----------------------------------------------------------------------------------------

\section{Einleitung}
    \begin{frame}{Einleitung}
    \begin{columns}
        \column{0.6\textwidth}
        %Begründung der Themenstellung aus dem derzeitigen Stand der Forschung: 
        %Warum ist das Thema relevant und was sind die offenen Fragen, die in der Präsentation behandelt werden sollen? An dieser Stelle soll auf die ausgewählte Fragestellung zurückgegriffen werden.

        \begin{itemize}
        \item \dots
        \end{itemize}

        \vspace*{0.75cm}

        %Ziel der Präsentation: Was soll dargestellt / herausgefunden werden und warum?

        \begin{itemize}
        \item \dots
        \end{itemize}        
        
        % Abgrenzung des Themas und evtl. themenbezogene Definitionen: Was kann in der Präsentation geleistet werden und was nicht? Welche Reichweite sollen die angestrebten Ergebnisse aufweisen?

        \begin{itemize}
            \item \dots
        \end{itemize}      

        % Überblick über Aufbau und Argumentationsfolge: Wie gliedert sich die Präsentation und was erwartet das Publikum im Hauptteil?

        \column{0.4\textwidth}

        
        % \begin{figure}
        %     \includegraphics[scale=1]{assets/bild.png}
        % \end{figure}
        % \vspace*{-0.5cm}
        % \scalebox{.4}{Schaubild (eigene Grafik)}

        \pnote{1-7}
            
        \pnote{1. \dots}
        
    \end{columns}
    \end{frame}

\section{Gängige Ansätze des semi-überwachten Lernens}
    \begin{frame}{Gängige Ansätze des semi-überwachten Lernens}

        % Aufgabenstellung 1.:
        %  Erkläre den Begriff des semi-überwachten Lernens und gib einen Überblick über gängige Ansätze, dieses Lernprinzip umzusetzen.

        \begin{columns}

            \column{0.6\textwidth}

                % \begin{figure}
                %     \includegraphics[scale=1.5]{assets/schaubild.png}
                % \end{figure}
                % \vspace*{-0.5cm}
                % \scalebox{.6}{Schaubild (Quelle: Wikimedia )}

            \column{0.4\textwidth}        

            \begin{itemize}
                \item \textbf{\dots:} \dots
            \end{itemize}


        \end{columns}
        
        \pnote{1-6}
        \pnote{1. \dots}
    \end{frame}

\section{Leistungsmessung bei semi-überwachten Lernverfahren}
    \begin{frame}{Leistungsmessung bei semi-überwachten Lernverfahren}

         %Aufgabenstellung 2.:
         % Beschreibe, wie sich die Performanz eines semi-überwachten Lernverfahrens messen lässt.

        \begin{alertblock}{\dots}
            \begin{itemize}
                \item \dots
            \end{itemize}
        \end{alertblock}
    
        \begin{exampleblock}{\dots}
            \begin{itemize}
                \item \dots
            \end{itemize}
        \end{exampleblock}
    
        \begin{extrablock}{\dots}
            \begin{itemize}
                \item \dots
            \end{itemize}
        \end{extrablock}
    
        \pnote{1-4}
        \pnote{1. ... }
      
    \end{frame}

\section{Szenarien für semi-überwachtes Lernen}

    \begin{frame}{Szenarien für semi-überwachtes Lernen}

        %Aufgabenstellung 3.:
        % Stelle Szenarien dar, in denen dieses semi-überwachte Lernen mit großer Wahrscheinlichkeit erfolgreich eingesetzt werden könnte.

        \begin{columns}
            \column{0.4\textwidth}

                % \begin{figure}
                %     \includegraphics[scale=0.17]{assets/symbolbild.jpeg}
                % \end{figure}
                % \vspace*{-0.5cm}
                % \scalebox{.4}{Symbolbild (generiert mit Bing Image Creator)}
            
                \column{0.6\textwidth}

            \begin{itemize}
                \item \textbf{\dots:} 
                \begin{itemize}
                    \item \dots
                \end{itemize}
                
                \item \textbf{\dots:} 
                \begin{itemize}
                    \item \dots
                \end{itemize}
                
                \item \textbf{\dots:} 
                \begin{itemize}
                    \item \dots
                \end{itemize}
                
                \item \textbf{\dots:} 
                \begin{itemize}
                    \item \dots
                \end{itemize}
            \end{itemize}

        \end{columns}

        \pnote{0-4}
        \pnote{0. ...}

    \end{frame}


\section{Zusammenfassung}

\begin{frame}{Zusammenfassung}

    \begin{itemize}
        \item \textbf{\dots:}
        \begin{itemize}
            \item \dots
        \end{itemize}
        
        \item \textbf{\dots:}
        \begin{itemize}
            \item \dots
        \end{itemize}
        
        \item \textbf{\dots:}
        \begin{itemize}
            \item \dots
        \end{itemize}
    \end{itemize}
    
    \pnote{1. ...}
    \pnote{8. ...}

    \end{frame}
    
\section{Referenzen}

\begin{frame}{Referenzen}
    
    \footnotesize{
    \begin{thebibliography}{99}

        \bibitem[IU, 2023]{p1} IU Internationale Hochschule GmbH (2023)
        \newblock Deep Learning - DLMDWDL01, Studienskript @IU International University of Applied Sciences

        \bibitem[Boyd, 2020]{p1} Umberto Michelucci (2019)
        \newblock Advanced Applied Deep Learning. Apress. \url{https://doi.org/10.1007/978-1-4842-4976-5}

    \end{thebibliography}
    }

    \pnote{1-2}
    \pnote{1. aus Zeitgründen nicht weiter drauf eingehen}
    \pnote{2. Danke für die Aufmerksamkeit und einen schönen Tag}
\end{frame}

\end{document} 
